\chapter*{ประวัติผู้เขียน}
\addcontentsline{toc}{chapter}{ประวัติผู้เขียน}

\vspace{0.4cm}
\begin{center}
    {\Large\bfseries ผศ.ดร.ชีวะ ทัศนา}\\[4pt]
    {\normalsize\bfseries สาขาวิชาฟิสิกส์ คณะวิทยาศาสตร์และเทคโนโลยี มหาวิทยาลัยราชภัฏรำไพพรรณี}\\[2pt]
    {\small อีเมล: \texttt{chewa.t@rbru.ac.th}}
\end{center}

\vspace{0.6cm}
\noindent
\begin{minipage}[t]{0.48\textwidth}
    \textbf{\large ประวัติการศึกษา}
    \begin{itemize}[leftmargin=1.2cm, itemsep=4pt]
        \item \textbf{ปร.ด. (ฟิสิกส์ประยุกต์)}\\
        มหาวิทยาลัยเทคโนโลยีสุรนารี
        \item \textbf{วท.ม. (ฟิสิกส์)}\\
        มหาวิทยาลัยเชียงใหม่
        \item \textbf{วท.บ. (ฟิสิกส์ เกียรตินิยม)}\\
        มหาวิทยาลัยเชียงใหม่
    \end{itemize}

    \vspace{0.4cm}
    \textbf{\large ประสบการณ์การสอนและบริหาร}
    \begin{itemize}[leftmargin=1.2cm, itemsep=4pt]
        \item ผู้ช่วยศาสตราจารย์ สาขาวิชาฟิสิกส์ มรภ.รำไพพรรณี
        \item อาจารย์ผู้รับผิดชอบหลักสูตรวิทยาศาสตรบัณฑิต
        \item ผู้ทรงคุณวุฒิประเมินบทความวิจัยระดับชาติและนานาชาติ
    \end{itemize}
\end{minipage}
\hfill
\begin{minipage}[t]{0.48\textwidth}
    \textbf{\large ความเชี่ยวชาญและงานวิจัย}
    \begin{itemize}[leftmargin=1.2cm, itemsep=4pt]
        \item \textbf{Computational Solid-State Physics:}\\
        การคำนวณโครงสร้างแถบพลังงาน $\text{LSDA}+U$ และอันตรกิริยาแลกเปลี่ยนแม่เหล็กในสารประกอบออกไซด์
        \item \textbf{Astrophysics \& Photometry:}\\
        การวิเคราะห์ระบบดาวคู่สัมผัสและการประมวลผลสัญญาณ CCD Photometry
        \item \textbf{Scientific Machine Learning:}\\
        การประยุกต์ใช้ PINNs และ Deep Learning ในการแก้สมการเชิงอนุพันธ์ย่อยทางฟิสิกส์
        \item \textbf{Environmental Physics:}\\
        การตรวจวัดอนุภาคละอองลอย PM2.5 และการประเมินพลังงานชีวมวล
    \end{itemize}
\end{minipage}
